%
% talk.tex
%
% Z Specification for Biff REPL Talk Subsystem
% Formal model of the two-phase handshake protocol, modal conversation,
% notification queue, and mutual hangup for the biff REPL talk command.
%

\documentclass[a4paper,10pt,fleqn]{article}
\usepackage[margin=1in]{geometry}
\usepackage{fuzz}
\usepackage[colorlinks=true,linkcolor=blue,citecolor=blue,urlcolor=blue]{hyperref}

\begin{document}

\title{Biff REPL Talk: A Z Specification}
\author{Formal Model of the Two-Phase Talk Handshake and Conversation Protocol}
\date{March 2026}
\hypersetup{
  pdftitle={Biff REPL Talk: A Z Specification},
  pdfauthor={punt-labs},
  pdfsubject={Z Specification},
  pdfcreator={fuzz/probcli}
}
\maketitle

\tableofcontents
\newpage

%%----------------------------------------------------------------------
\section{Introduction}
\label{sec:introduction}
%%----------------------------------------------------------------------

The biff REPL talk command implements a modal conversation between two
sessions, modelled on BSD \texttt{talk(1)}.  The protocol has three
phases:

\begin{enumerate}
\item \textbf{Handshake} --- the inviter publishes an \textit{invite}
  notification and enters a waiting state.  The responder sees the
  invite as a banner, types \texttt{talk @inviter:tty}, and publishes
  an \textit{accept} notification.  The inviter receives the accept
  and both sides enter the connected state.

\item \textbf{Conversation} --- both sides exchange \textit{message}
  notifications via NATS core publish (ephemeral, no inbox).  Messages
  are received via a notification queue drained on a 2-second timeout
  poll.

\item \textbf{Hangup} --- either side types \texttt{end}, which
  publishes an \textit{end} notification.  The remote side detects it
  during the drain and exits talk mode.  Both sides return to the REPL.
\end{enumerate}

This specification models the state of a single session's talk
subsystem and the operations that transition it.  A two-session
system is modelled by composing two instances.

%%----------------------------------------------------------------------
\section{Basic Types}
%%----------------------------------------------------------------------

\begin{zed}
  [USER, TTY, MSGBODY, SESSIONKEY]
\end{zed}

\begin{zed}
  [ORG, REPO]
\end{zed}

$USER$ identifies a biff user.  $TTY$ identifies a session (hex ID or
tty name).  $MSGBODY$ is an opaque message body (max 512 chars).
$SESSIONKEY$ is a composite \texttt{user:tty} key --- the globally-unique
talk address (DES-035).  $ORG$ is an organization and $REPO$ a
repository.  A session runs in exactly one repository, which belongs to
exactly one organization --- the containment hierarchy of
$session\-model.tex$, imported here through the resolver of
Section~\ref{sec:addressing}.

A \emph{subject} is the NATS topic a talk frame is published to.  Talk
routes on the globally-unique identity \emph{alone}: the subject is the
session key, and nothing else.  Neither the organization nor the
repository appears in it --- both are visibility and presence attributes
of the session, not coordinates of the route.  In the running system
the subject is \texttt{\{prefix\}.talk.notify.\{user\}:\{tty\}}.

\begin{zed}
  SUBJECT == SESSIONKEY
\end{zed}

%%----------------------------------------------------------------------
\section{Free Types}
%%----------------------------------------------------------------------

\begin{zed}
  ZBOOL ::= ztrue | zfalse
\end{zed}

\begin{zed}
  TalkPhase ::= tpIdle | tpInviting | tpConnected
\end{zed}

The three phases of a talk session from one side's perspective:
\begin{itemize}
\item $tpIdle$ --- normal REPL mode, not in talk.
\item $tpInviting$ --- sent an invite, waiting for accept.
\item $tpConnected$ --- in the talk sub-loop, exchanging messages.
\end{itemize}

Note: there is no explicit ``accepting'' phase.  When a responder
types \texttt{talk @inviter:tty} and a pending invite exists, the
accept is sent and the state transitions directly from $tpIdle$
to $tpConnected$ (via $RespondToInvite$).

\begin{zed}
  NotifType ::= ntInvite | ntAccept | ntMessage | ntEnd
\end{zed}

The four notification types in the talk protocol.  These correspond
to the \texttt{type} field in the NATS JSON payload.

\begin{zed}
  ROUTEMODE ::= rmIdentity | rmRepoScoped | rmOrgScoped
\end{zed}

$ROUTEMODE$ selects the routing discipline a session uses when it
publishes a reply.  Under $rmIdentity$ --- the correct design --- the
reply is addressed to the peer's globally-unique identity, taken from the
frame it is replying to; nothing about the peer beyond its identity is
consulted, and any resolvable peer is reached whatever its organization
or repository.  The other two modes are the over-scoped defects.  Under
$rmRepoScoped$ the subject is keyed on the repository and the responder
routes on its \emph{own} repository, so a reply to a peer in another
repository lands on the wrong subject.  Under $rmOrgScoped$ the subject
is keyed on the organization and the responder routes on its \emph{own}
organization, so a reply to a peer in another organization lands on the
wrong subject.  The mode is fixed at initialization and never changes; it
is a scoping parameter that lets one model exhibit the correct design
against \emph{both} over-scoped defects (Section~\ref{sec:routing}).

%%----------------------------------------------------------------------
\section{Addressing and Routing}
\label{sec:addressing}
%%----------------------------------------------------------------------

Talk names its peer by the globally-unique session identity
\texttt{@user:tty} (DES-035): tty names are reserved atomically across
all repositories, so a $SESSIONKEY$ denotes one session regardless of
which repository or organization it runs in.  Routing is a function of
that identity and of nothing else.  The organization and the repository
are \emph{attributes} of the addressed session, used for presence and
visibility; neither is a coordinate of the route.

The resolver below is the interface to $session\-model.tex$.  That
specification owns the entities (organizations, repositories, members,
sessions, the $srepo$ attribute); this one owns only the fact that every
live session key resolves to a user, a repository, and an organization,
and that its subject is its identity.  We do not restate $srepo$ here;
$keyRepo$ and $keyOrg$ \emph{are} $srepo$ viewed through the
\texttt{@user:tty} address, and both are visibility attributes, not
routing coordinates (see Section~\ref{sec:ownership}).

\begin{axdef}
  keyUser : SESSIONKEY \fun USER \\
  keyRepo : SESSIONKEY \surj REPO \\
  repoOrg : REPO \surj ORG \\
  keyOrg : SESSIONKEY \fun ORG \\
  subjectOf : SESSIONKEY \fun SUBJECT
\where
  \forall k : SESSIONKEY @ keyOrg~k = repoOrg~(keyRepo~k) \\
  \forall k : SESSIONKEY @ subjectOf~k = k
\end{axdef}

$keyUser~k$, $keyRepo~k$ and $keyOrg~k$ are the user, repository and
organization of the session $k$; $repoOrg$ carries each repository to its
owning organization.  $subjectOf~k = k$ is the subject that session $k$
subscribes to and that any correctly routed frame for $k$ must carry: the
identity is the whole of the route.  A sender routing a reply to peer $p$
needs only $p$'s identity, taken from the frame it is replying to; it
resolves neither $p$'s organization nor $p$'s repository, and persists
nothing.  $keyRepo$ and $keyOrg$ survive in the resolver only to state
that $srepo$ and its organization are presence attributes, and to give
the two \emph{defective} disciplines a coordinate to mis-route on:
$rmRepoScoped$ mis-routes on $keyRepo$, $rmOrgScoped$ on $keyOrg$.

$repoOrg$ is surjective ($\surj$) as well as $keyRepo$: every modelled
organization owns a repository, and every repository holds a session.  At
$\#SESSIONKEY = 2$, $\#REPO = 2$, $\#ORG = 2$ this forces the two keys
into \emph{different} repositories \emph{and} different organizations, so
the peer a session addresses differs from it on both attributes at once
--- the configuration under which \emph{both} over-scoped defects
mis-route, and under which pure-identity routing must still deliver.

Surjectivity is not a loss of generality --- the repositories are
exactly those org discovery enumerates (DES-034 lists only repositories
\emph{with active sessions}), and each organization in play owns one ---
and it has a sharp verification purpose.  It removes the constant
solver's discretion: at two keys, two repositories and two
organizations the peer is guaranteed to differ from the sender in both
repository and organization, so a single scope exercises the repo-keyed
strand, the org-keyed strand, and cross-attribute delivery under
pure-identity routing at once.  At $\#REPO = \#ORG = 1$ the model
degenerates to a single repository and organization, in which no reply
can be mis-routed by either defect.

%%----------------------------------------------------------------------
\section{Global Constants}
%%----------------------------------------------------------------------

\begin{zed}
  [QSLOT]
\end{zed}

\begin{axdef}
  maxQueueLen : \nat_1 \\
  maxBodyLen : \nat
\where
  QSLOT \neq \emptyset \\
  maxQueueLen = \# QSLOT \\
  maxBodyLen = 512
\end{axdef}

$maxQueueLen$ bounds the notification queue.  In the running system the
queue is capped at a fixed $MAX\_TALK\_QUEUE = 100$ with drop-oldest
eviction (DES-044).  In this model $maxQueueLen$ is deliberately
\emph{not} pinned to $100$: it is the cardinality of an abstract carrier
$QSLOT$, the set of available queue slots, so the bound is a
\emph{scoping parameter} ProB can shrink for exhaustive checking rather
than a hard-wired value.  $QSLOT$ is required non-empty, so
$maxQueueLen \geq 1$ --- a queue with at least one slot, which keeps
$ReceiveOverflow$ well defined ($tail$ is applied only to a non-empty
queue).  Because $QSLOT$ is a given set, its size is chosen by ProB ---
uniformly by \texttt{DEFAULT\_SETSIZE}, or independently of the other
carriers by \texttt{-card QSLOT}~$n$.  This lets an exhaustive check
bound the queue small while keeping two distinct session keys, which is
what exercises $MutualAutoAccept$.  The specification asserts that the
queue is bounded by some fixed natural number $\geq 1$; it deliberately
does not fix that number to the concrete $100$ the implementation uses
(DES-044).  $maxBodyLen$ is the message body truncation limit.

The model is verified at two complementary scopes.  No single small
scope covers every operation at once: the routing operations need two
repositories and two organizations, and the two forged connected-frame
operations need a third session key.  Between them the two runs cover
every operation, and each holds every invariant with no deadlock.  The
addressing layer adds the constant resolver ($keyUser$, $keyRepo$,
$repoOrg$) whose valuations the constant solver enumerates, and the
$visible$ set adds a per-run peer configuration, so a full run at
$\texttt{DEFAULT\_SETSIZE}$ 2 (all carriers 2) explores tens of
thousands of states; the targeted scopes below are exhaustive over the
behaviour that matters and stay small.  $MAX\_INITIALISATIONS$ must be
set generously --- $Init$ leaves $routeMode$, $visible$ and the identity
free, and a value too small silently truncates the initial states (ProB
warns \texttt{not all transitions were computed}); $512$ suffices here.

\emph{Routing run} --- two session keys in two repositories of two
organizations.  Surjective $keyRepo$ and $repoOrg$ put the two keys in
different repositories \emph{and} different organizations, so the peer a
session addresses differs from it on both attributes at once; the two
keys and the $keyBelow$ order supply the mutual-invite tie-break, and the
free $visible$ set supplies the peer configurations:

\begin{verbatim}
probcli docs/talk.tex -model_check -noltl -coverage \
  -card USER 1 -card TTY 1 -card MSGBODY 1 \
  -card SESSIONKEY 2 -card REPO 2 -card ORG 2 -card QSLOT 1 \
  -p MAX_INITIALISATIONS 512 -p MAX_OPERATIONS 12000
\end{verbatim}

All $2{,}089$ reachable states are visited with no invariant violation
and no deadlock, and every operation is covered save $DrainForeignMessage$
and $DrainForeignEnd$ (below).  In particular the delivery-soundness
invariant $routeMode = rmIdentity \implies unsoundEmit = zfalse$ holds
throughout, $ReceiveNotForSubject$ fires (a frame on a foreign subject is
dropped), and $MutualAutoAccept$ is exercised.  The routing properties ---
soundness, each defect's strand, and unreachability under identity
routing --- are checked as reachability goals against this same scope
(Section~\ref{sec:routing}).

\emph{Forgery run} --- three session keys in one repository and
organization.  The two connected-phase foreign-frame operations,
$DrainForeignMessage$ and $DrainForeignEnd$, guard on $phase =
tpConnected$ with $(head~queue).nfromKey \neq partnerKey$: a three-party
predicate needing $myKey$, the connected $partnerKey$, and a third
\emph{forger} key distinct from both.  Two keys cannot supply the third,
so these operations are enabled in no reachable state at the routing
scope; exhausting them takes a third key.  Routing is irrelevant to the
threat, so a single repository and organization keeps the check small:

\begin{verbatim}
probcli docs/talk.tex -model_check -noltl -coverage \
  -card USER 1 -card TTY 1 -card MSGBODY 1 \
  -card SESSIONKEY 3 -card REPO 1 -card ORG 1 -card QSLOT 1 \
  -p MAX_INITIALISATIONS 512 -p MAX_OPERATIONS 20000
\end{verbatim}

All $9{,}727$ reachable states are visited, no invariant violation and
no deadlock, and \emph{every} operation is covered --- the third key
supplies both the forger the $DrainForeign$ guards need and, since the
subject keys on identity, a foreign subject for $ReceiveNotForSubject$
to drop even at one repository.  A forged message or end carrying the
third key is drained by $DrainForeignMessage$ or $DrainForeignEnd$ and
never reaches $DrainMessage$ or $DrainEnd$, so it neither speaks in the
partner's place nor hangs up the call.  Consent still keys on the
ephemeral session key, exactly as before the addressing change: the
subject decides \emph{where} a frame goes, the session key decides
\emph{whether} it is honoured, and the two are orthogonal (DES-046).
$MAX\_INITIALISATIONS$ must exceed the number of identities $Init$ leaves
unconstrained, or ProB silently drops initial states; at three keys with
the two $keyBelow$ orders, $128$ suffices.

$keyBelow$ is a strict total order on session keys, modelling the
lexicographic comparison the implementation uses to break mutual-invite
ties deterministically.  Because it is total and antisymmetric, for any
two distinct keys exactly one is below the other --- so exactly one side
of a simultaneous mutual invite auto-accepts and neither blocks forever.

\begin{axdef}
  keyBelow : SESSIONKEY \rel SESSIONKEY
\where
  \forall k : SESSIONKEY @ (k \mapsto k) \notin keyBelow \\
  \forall j, k : SESSIONKEY @
    (j \mapsto k) \in keyBelow \implies (k \mapsto j) \notin keyBelow \\
  \forall j, k, l : SESSIONKEY @
    (j \mapsto k) \in keyBelow \land (k \mapsto l) \in keyBelow
      \implies (j \mapsto l) \in keyBelow \\
  \forall j, k : SESSIONKEY @
    j \neq k \implies (j \mapsto k) \in keyBelow \lor (k \mapsto j) \in keyBelow
\end{axdef}

%%----------------------------------------------------------------------
\section{Notification}
%%----------------------------------------------------------------------

A notification is a typed message in the queue, carrying the sender
identity and an optional body.

\begin{schema}{Notification}
  ntype : NotifType \\
  nfrom : USER \\
  nfromTty : TTY \\
  nfromKey : SESSIONKEY \\
  nto : SESSIONKEY \\
  nbody : MSGBODY
\end{schema}

$nto$ is the target session key.  Talk is session-scoped
(DES-043): every notification names exactly one recipient session.
The per-identity NATS subject carries the notification to the addressed
identity's session (Section~\ref{sec:addressing}), and $nto$ is the
session-scoping check on top --- a session accepts a notification only
when $nto$ matches its own key.

%%----------------------------------------------------------------------
\section{Talk State}
%%----------------------------------------------------------------------

The talk state for a single REPL session.

\begin{schema}{TalkState}
  phase : TalkPhase \\
  myUser : USER \\
  myTty : TTY \\
  myKey : SESSIONKEY \\
  myRepo : REPO \\
  myOrg : ORG \\
  partner : USER \\
  partnerTty : TTY \\
  partnerKey : SESSIONKEY \\
  pendingInvites : USER \pfun SESSIONKEY \\
  queue : \seq Notification \\
  queueLen : \nat \\
  visible : \power SESSIONKEY \\
  routeMode : ROUTEMODE \\
  unsoundEmit : ZBOOL
\where
  queueLen = \# queue \\
  queueLen \leq maxQueueLen \\
  phase = tpIdle \implies partner = myUser \\
  phase = tpIdle \implies partnerTty = myTty \\
  phase = tpIdle \implies partnerKey = myKey \\
  myUser = keyUser~myKey \\
  myRepo = keyRepo~myKey \\
  myOrg = keyOrg~myKey \\
  myKey \notin visible \\
  (phase \neq tpIdle \implies partnerKey \in visible) \\
  routeMode = rmIdentity \implies unsoundEmit = zfalse
\end{schema}

When idle, $partner$, $partnerTty$, and $partnerKey$ are set to
$myUser$, $myTty$, and $myKey$ as sentinel values (no active partner).
$partnerKey$ records the exact session key we invited (or that invited
us); it is the consent anchor for the accept handshake --- an accept is
honoured only when its origin equals $partnerKey$ (DES-043).  The
$pendingInvites$ partial function maps each user who has sent us an
invite to that invite's originating session key --- the key a
subsequent \texttt{talk} command uses to target its accept.  Because
it is a function, a newer invite from the same user \emph{supersedes}
the older one (the maplet is overwritten).  $queue$ is the
notification queue populated by the NATS subscription callback.

Here $pendingInvites$ is keyed by $USER$, a finite given set, so its
size is intrinsically bounded and this model does not cap it
explicitly.  The running system is not so lucky: it keys the pending
set by the inviter's \emph{wire-controlled} identity, an unbounded
string, so a stream of forged inviters would grow it without limit ---
the same unbounded-growth denial of service the queue cap $maxQueueLen$
(DES-044) guards against.  The sister specification $notification.tex$
therefore bounds the pending set by $maxPending$ ($\# PSLOT$) and evicts
the oldest entry on overflow ($TalkInviteArriveOverflow$), the
pending-set mirror of this queue's $maxQueueLen$ and $ReceiveOverflow$.

The last five components carry the addressing layer.  $myRepo$ and
$myOrg$ are the repository and organization of this session; they are
not free, being pinned to the identity by $myRepo = keyRepo~myKey$ and
$myOrg = keyOrg~myKey$, and they are here only as presence attributes and
as the coordinates the two defective disciplines wrongly route on.  The
subscription subject is $subjectOf~myKey = myKey$ --- the identity, and
nothing else.  $visible$ is the set of sessions this session can resolve
--- its \emph{peers}, as fixed by the \texttt{peers} and \texttt{orgs}
configuration.  It excludes $myKey$ ($myKey \notin visible$) and is the
only gate on talk: an invite or an accept may name a peer only if that
peer is in $visible$ (Sections~\ref{sec:addressing},~\ref{sec:routing}).
The invariant $phase \neq tpIdle \implies partnerKey \in visible$ records
that a partner is always someone we can see, which is what makes it
addressable at all.  Visibility spans organizations and repositories
freely: a visible peer is reachable whatever its $keyOrg$ or $keyRepo$.
$routeMode$ fixes the routing discipline for the whole run.  $unsoundEmit$
is a history flag: it is $zfalse$ initially and latches to $ztrue$ the
first time this session publishes a reply to a subject other than its
addressee's own ($subjectOf$ of the key it is replying to).  The state
invariant $routeMode = rmIdentity \implies unsoundEmit = zfalse$ is the
delivery-soundness guarantee: under identity routing no session ever
mis-routes a reply, so every reply reaches the addressed
\texttt{@user:tty} whichever organization or repository either party is
in.  That the same flag \emph{can} latch under $rmRepoScoped$ or
$rmOrgScoped$ is the over-scoping defect, reproduced in
Section~\ref{sec:routing}.

%%----------------------------------------------------------------------
\section{Initialization}
%%----------------------------------------------------------------------

\begin{schema}{Init}
  TalkState'
\where
  phase' = tpIdle \\
  pendingInvites' = \emptyset \\
  queue' = \langle \rangle \\
  queueLen' = 0 \\
  partner' = myUser' \\
  partnerTty' = myTty' \\
  partnerKey' = myKey' \\
  unsoundEmit' = zfalse
\end{schema}

$routeMode'$ and $visible'$ are left unconstrained: $Init$ admits every
routing discipline and every peer-set the invariants allow, and the
operations never reassign either, so ProB explores each discipline and
each visibility configuration as a separate, constant branch.  $myRepo'$
and $myOrg'$ are pinned by the $TalkState$ invariant to the initial
identity, and $visible'$ excludes $myKey'$.

%%----------------------------------------------------------------------
\section{Operations}
%%----------------------------------------------------------------------

%%----------------------------------------------------------------------
\subsection{ReceiveNotification}
%%----------------------------------------------------------------------

A notification arrives from NATS and is appended to the queue.
Self-echo (notifications from our own session key) is silently
dropped.  A notification whose target session key $nto$ is not our
own key is also dropped --- talk is session-scoped, so only the named
session accepts it (DES-043).  The queue is bounded; on overflow the
oldest queued notification is dropped to make room (ReceiveOverflow).

\begin{schema}{ReceiveNotification}
  \Delta TalkState \\
  notif? : Notification \\
  insubject? : SUBJECT
\where
  notif?.nfromKey \neq myKey \\
  notif?.nto = myKey \\
  insubject? = myKey \\
  queueLen < maxQueueLen \\
  queue' = queue \cat \langle notif? \rangle \\
  queueLen' = queueLen + 1 \\
  phase' = phase \\
  myUser' = myUser \\
  myTty' = myTty \\
  myKey' = myKey \\
  myRepo' = myRepo \\
  myOrg' = myOrg \\
  partner' = partner \\
  partnerTty' = partnerTty \\
  partnerKey' = partnerKey \\
  pendingInvites' = pendingInvites \\
  visible' = visible \\
  routeMode' = routeMode \\
  unsoundEmit' = unsoundEmit
\end{schema}

The new input $insubject?$ is the NATS subject the frame arrived on.  A
session subscribes to exactly its own subject $subjectOf~myKey = myKey$;
a frame is delivered to it only when it was published there.  This makes
the transport explicit: identity routing happens before the $nto$
session-scoping filter, not after it.  There is no organization guard on
receipt: talk is not org-scoped, and a frame from a session in any
organization may arrive --- this is also the honest model of the forgeable
transport, on which a frame from any key at all may be published to our
subject (the threat model of Section~\ref{sec:threat}).  A frame carrying
the correct $nto$ but published to a foreign subject never arrives ---
$ReceiveNotForSubject$ below --- which is the receiving half of the
delivery argument (Section~\ref{sec:routing}).

%%----------------------------------------------------------------------
\subsection{ReceiveSelfEcho}
%%----------------------------------------------------------------------

A notification from our own session key is silently dropped.

\begin{schema}{ReceiveSelfEcho}
  \Xi TalkState \\
  notif? : Notification
\where
  notif?.nfromKey = myKey
\end{schema}

%%----------------------------------------------------------------------
\subsection{ReceiveNotForSession}
%%----------------------------------------------------------------------

A notification addressed to a different session ($nto \neq myKey$) is
silently dropped.  This is the session-scoping filter: the per-user
subject is shared, but only the named session accepts.

\begin{schema}{ReceiveNotForSession}
  \Xi TalkState \\
  notif? : Notification
\where
  notif?.nfromKey \neq myKey \\
  notif?.nto \neq myKey
\end{schema}

%%----------------------------------------------------------------------
\subsection{ReceiveNotForSubject}
%%----------------------------------------------------------------------

A frame published to a subject other than our own never reaches us: we
subscribe to $myKey$ alone, and the transport delivers on a subject only
to its subscribers.  This is the receiving half of the routing argument.
A frame may carry a correct $nto = myKey$ and still be lost here, if its
publisher addressed the wrong subject --- which is exactly what a
repo-scoped or org-scoped reply to a peer in another repository or
organization does (Section~\ref{sec:routing}).  The frame is dropped with
no change of state.

\begin{schema}{ReceiveNotForSubject}
  \Xi TalkState \\
  notif? : Notification \\
  insubject? : SUBJECT
\where
  notif?.nfromKey \neq myKey \\
  insubject? \neq myKey
\end{schema}

%%----------------------------------------------------------------------
\subsection{ReceiveOverflow}
%%----------------------------------------------------------------------

Queue is full; the oldest queued notification is dropped and the new
one appended (drop-oldest).  The newest $maxQueueLen$ notifications are
retained; length stays at $maxQueueLen$.

\begin{schema}{ReceiveOverflow}
  \Delta TalkState \\
  notif? : Notification \\
  insubject? : SUBJECT
\where
  queueLen = maxQueueLen \\
  notif?.nfromKey \neq myKey \\
  notif?.nto = myKey \\
  insubject? = myKey \\
  queue' = tail~queue \cat \langle notif? \rangle \\
  queueLen' = queueLen \\
  phase' = phase \\
  myUser' = myUser \\
  myTty' = myTty \\
  myKey' = myKey \\
  myRepo' = myRepo \\
  myOrg' = myOrg \\
  partner' = partner \\
  partnerTty' = partnerTty \\
  partnerKey' = partnerKey \\
  pendingInvites' = pendingInvites \\
  visible' = visible \\
  routeMode' = routeMode \\
  unsoundEmit' = unsoundEmit
\end{schema}

%%----------------------------------------------------------------------
\subsection{DrainInvite}
%%----------------------------------------------------------------------

During the REPL notification poll, an invite is drained from the
queue.  The sender is mapped to the invite's session key in
$pendingInvites$ so that a subsequent \texttt{talk} command can target
its accept at that exact session.  Function override ($\oplus$) means a
newer invite from the same user supersedes the older one.

\begin{schema}{DrainInvite}
  \Delta TalkState
\where
  phase = tpIdle \\
  queue \neq \langle \rangle \\
  (head~queue).ntype = ntInvite \\
  pendingInvites' =
    pendingInvites \oplus \{ (head~queue).nfrom \mapsto (head~queue).nfromKey \} \\
  queue' = tail~queue \\
  queueLen' = queueLen - 1 \\
  phase' = phase \\
  myUser' = myUser \\
  myTty' = myTty \\
  myKey' = myKey \\
  myRepo' = myRepo \\
  myOrg' = myOrg \\
  partner' = partner \\
  partnerTty' = partnerTty \\
  partnerKey' = partnerKey \\
  visible' = visible \\
  routeMode' = routeMode \\
  unsoundEmit' = unsoundEmit
\end{schema}

%%----------------------------------------------------------------------
\subsection{DrainAcceptWhileInviting}
%%----------------------------------------------------------------------

While in the inviting phase, an accept from the very session we invited
is drained from the queue.  This completes the handshake and transitions
to connected.  The consent guard $(head~queue).nfromKey = partnerKey$ is
essential: only the invited session's accept may connect us.  An accept
from any other origin --- even one correctly addressed to us
($nto = myKey$) --- must not complete the handshake, or a third party
could forge the invited peer's consent (DES-043).

\begin{schema}{DrainAcceptWhileInviting}
  \Delta TalkState
\where
  phase = tpInviting \\
  queue \neq \langle \rangle \\
  (head~queue).ntype = ntAccept \\
  (head~queue).nfromKey = partnerKey \\
  phase' = tpConnected \\
  queue' = tail~queue \\
  queueLen' = queueLen - 1 \\
  myUser' = myUser \\
  myTty' = myTty \\
  myKey' = myKey \\
  myRepo' = myRepo \\
  myOrg' = myOrg \\
  partner' = partner \\
  partnerTty' = partnerTty \\
  partnerKey' = partnerKey \\
  pendingInvites' = pendingInvites \\
  visible' = visible \\
  routeMode' = routeMode \\
  unsoundEmit' = unsoundEmit
\end{schema}

%%----------------------------------------------------------------------
\subsection{DrainForeignAccept}
%%----------------------------------------------------------------------

An accept whose origin is not the invited session ($nfromKey \neq
partnerKey$) is drained and discarded without changing phase.  This is
the model counterpart of the consent guard above: the initiator keeps
waiting.  It also covers a stray accept arriving in any phase, matching
the implementation, which drains the whole queue and honours only the
invited peer's accept.

\begin{schema}{DrainForeignAccept}
  \Delta TalkState
\where
  queue \neq \langle \rangle \\
  (head~queue).ntype = ntAccept \\
  (head~queue).nfromKey \neq partnerKey \\
  queue' = tail~queue \\
  queueLen' = queueLen - 1 \\
  phase' = phase \\
  myUser' = myUser \\
  myTty' = myTty \\
  myKey' = myKey \\
  myRepo' = myRepo \\
  myOrg' = myOrg \\
  partner' = partner \\
  partnerTty' = partnerTty \\
  partnerKey' = partnerKey \\
  pendingInvites' = pendingInvites \\
  visible' = visible \\
  routeMode' = routeMode \\
  unsoundEmit' = unsoundEmit
\end{schema}

%%----------------------------------------------------------------------
\subsection{DrainMessage}
%%----------------------------------------------------------------------

During talk mode, a message notification \emph{from the connected
partner} is drained and displayed to the user.  The consent guard
$(head~queue).nfromKey = partnerKey$ binds the frame to the invited
peer: only the partner may speak.  A message whose origin is any other
session --- even one correctly addressed to us ($nto = myKey$) --- must
not be surfaced in the partner's place, or a third party could
impersonate the peer during a live conversation (DES-043).  Such a
frame is drained without display by $DrainForeignMessage$.

\begin{schema}{DrainMessage}
  \Delta TalkState
\where
  phase = tpConnected \\
  queue \neq \langle \rangle \\
  (head~queue).ntype = ntMessage \\
  (head~queue).nfromKey = partnerKey \\
  queue' = tail~queue \\
  queueLen' = queueLen - 1 \\
  phase' = phase \\
  myUser' = myUser \\
  myTty' = myTty \\
  myKey' = myKey \\
  myRepo' = myRepo \\
  myOrg' = myOrg \\
  partner' = partner \\
  partnerTty' = partnerTty \\
  partnerKey' = partnerKey \\
  pendingInvites' = pendingInvites \\
  visible' = visible \\
  routeMode' = routeMode \\
  unsoundEmit' = unsoundEmit
\end{schema}

%%----------------------------------------------------------------------
\subsection{DrainForeignMessage}
%%----------------------------------------------------------------------

A connected-phase message whose origin is not the connected partner
($nfromKey \neq partnerKey$) is drained and discarded without display.
This is the conversation-side counterpart of $DrainForeignAccept$: the
frame is dequeued so the poll makes progress, but it neither reaches
the user nor changes any component beyond the queue.  A forged
$ntMessage$ addressed to us ($nto = myKey$) is thereby prevented from
speaking in the partner's place.

\begin{schema}{DrainForeignMessage}
  \Delta TalkState
\where
  phase = tpConnected \\
  queue \neq \langle \rangle \\
  (head~queue).ntype = ntMessage \\
  (head~queue).nfromKey \neq partnerKey \\
  queue' = tail~queue \\
  queueLen' = queueLen - 1 \\
  phase' = phase \\
  myUser' = myUser \\
  myTty' = myTty \\
  myKey' = myKey \\
  myRepo' = myRepo \\
  myOrg' = myOrg \\
  partner' = partner \\
  partnerTty' = partnerTty \\
  partnerKey' = partnerKey \\
  pendingInvites' = pendingInvites \\
  visible' = visible \\
  routeMode' = routeMode \\
  unsoundEmit' = unsoundEmit
\end{schema}

%%----------------------------------------------------------------------
\subsection{DrainEnd}
%%----------------------------------------------------------------------

During talk mode, an end notification \emph{from the connected
partner} is drained: the peer has hung up, and the session transitions
back to idle.  The consent guard $(head~queue).nfromKey = partnerKey$
is essential --- only the connected partner may end the call.  A forged
$ntEnd$ from any other origin must not tear down the conversation, or a
third party could hang up a call it is not part of (DES-043); such a
frame is drained by $DrainForeignEnd$.

\begin{schema}{DrainEnd}
  \Delta TalkState
\where
  phase = tpConnected \\
  queue \neq \langle \rangle \\
  (head~queue).ntype = ntEnd \\
  (head~queue).nfromKey = partnerKey \\
  phase' = tpIdle \\
  partner' = myUser \\
  partnerTty' = myTty \\
  partnerKey' = myKey \\
  queue' = tail~queue \\
  queueLen' = queueLen - 1 \\
  myUser' = myUser \\
  myTty' = myTty \\
  myKey' = myKey \\
  myRepo' = myRepo \\
  myOrg' = myOrg \\
  pendingInvites' = pendingInvites \\
  visible' = visible \\
  routeMode' = routeMode \\
  unsoundEmit' = unsoundEmit
\end{schema}

%%----------------------------------------------------------------------
\subsection{DrainForeignEnd}
%%----------------------------------------------------------------------

A connected-phase end frame whose origin is not the connected partner
($nfromKey \neq partnerKey$) is drained and discarded, leaving the
phase unchanged.  This is the hangup-side counterpart of
$DrainForeignAccept$: the call continues, and only the partner's
$ntEnd$ can end it.  A forged $ntEnd$ addressed to us ($nto = myKey$)
is thereby prevented from tearing down a conversation it is not part
of.

\begin{schema}{DrainForeignEnd}
  \Delta TalkState
\where
  phase = tpConnected \\
  queue \neq \langle \rangle \\
  (head~queue).ntype = ntEnd \\
  (head~queue).nfromKey \neq partnerKey \\
  queue' = tail~queue \\
  queueLen' = queueLen - 1 \\
  phase' = phase \\
  myUser' = myUser \\
  myTty' = myTty \\
  myKey' = myKey \\
  myRepo' = myRepo \\
  myOrg' = myOrg \\
  partner' = partner \\
  partnerTty' = partnerTty \\
  partnerKey' = partnerKey \\
  pendingInvites' = pendingInvites \\
  visible' = visible \\
  routeMode' = routeMode \\
  unsoundEmit' = unsoundEmit
\end{schema}

%%----------------------------------------------------------------------
\subsection{SendInvite}
%%----------------------------------------------------------------------

User types \texttt{talk @target:tty} with no pending invite from
target.  Publishes a session-scoped invite (whose $nto$ is
$targetKey?$) and enters the inviting phase.  Talk is session-scoped,
so the address must name a specific session; self-talk is disallowed
($targetKey? \neq myKey$).  The one gate on the address is
\emph{visibility}: $targetKey? \in visible$.  You may invite any session
you can see and no session you cannot; a visible peer is reachable
whatever its organization or repository, and an invisible one is not
addressable at all.  This is the whole of the routing precondition ---
there is no organization or repository constraint on top of it.

The invite is the \emph{sound} direction, and the model reflects that:
the inviter addresses the target's own identity subject $subjectOf~
targetKey? = targetKey?$, so the invite reaches it under every
discipline.  The invite therefore never latches $unsoundEmit$
($unsoundEmit' = unsoundEmit$).  The defect lives entirely in the
\emph{reply} direction --- $RespondToInvite$, $MutualAutoAccept$,
$SendMessage$, $LocalEnd$ --- where an over-scoped design forces the
responder to route on its own repository or organization instead of the
peer's identity.  The correct design has no such asymmetry: a reply
routes to the peer's identity, held locally, so there is nothing to
resolve and nothing to persist.

\begin{schema}{SendInvite}
  \Delta TalkState \\
  target? : USER \\
  targetTty? : TTY \\
  targetKey? : SESSIONKEY
\where
  phase = tpIdle \\
  target? \notin \dom pendingInvites \\
  targetKey? \neq myKey \\
  targetKey? \in visible \\
  phase' = tpInviting \\
  partner' = target? \\
  partnerTty' = targetTty? \\
  partnerKey' = targetKey? \\
  queue' = queue \\
  queueLen' = queueLen \\
  myUser' = myUser \\
  myTty' = myTty \\
  myKey' = myKey \\
  myRepo' = myRepo \\
  myOrg' = myOrg \\
  pendingInvites' = pendingInvites \\
  visible' = visible \\
  routeMode' = routeMode \\
  unsoundEmit' = unsoundEmit
\end{schema}

%%----------------------------------------------------------------------
\subsection{RespondToInvite}
%%----------------------------------------------------------------------

User types \texttt{talk @inviter:tty} and a pending invite from
that user exists.  Consumes the invite, publishes an accept, and
transitions directly to connected.

\begin{schema}{RespondToInvite}
  \Delta TalkState \\
  target? : USER \\
  targetTty? : TTY
\where
  phase = tpIdle \\
  target? \in \dom pendingInvites \\
  pendingInvites~target? \in visible \\
  phase' = tpConnected \\
  partner' = target? \\
  partnerTty' = targetTty? \\
  partnerKey' = pendingInvites~target? \\
  pendingInvites' = \{ target? \} \ndres pendingInvites \\
  queue' = queue \\
  queueLen' = queueLen \\
  myUser' = myUser \\
  myTty' = myTty \\
  myKey' = myKey \\
  myRepo' = myRepo \\
  myOrg' = myOrg \\
  visible' = visible \\
  routeMode' = routeMode \\
  ((unsoundEmit' = ztrue) \iff
    (unsoundEmit = ztrue \lor
      (routeMode = rmRepoScoped \land
        myRepo \neq keyRepo~(pendingInvites~target?)) \lor
      (routeMode = rmOrgScoped \land
        myOrg \neq keyOrg~(pendingInvites~target?))))
\end{schema}

You may accept only an inviter you can resolve: $pendingInvites~target?
\in visible$.  This is the same visibility gate as $SendInvite$ --- an
invite from a session you cannot see is left in $pendingInvites$
unanswerable --- and it is what keeps the partner visible, hence the
non-idle invariant $partnerKey \in visible$.

The accept this operation publishes is addressed to the inviter, the
session key $pendingInvites~target?$.  This is the frame today's defect
mis-routes.  Under the correct design ($rmIdentity$) it is stamped
$subjectOf~(pendingInvites~target?) = pendingInvites~target?$ --- the
inviter's own identity, taken from the invite --- so the accept lands
where the inviter listens, and the acceptor needs no knowledge of the
inviter's organization or repository.  Under a defect the subject is
keyed on an attribute the acceptor reads from \emph{itself}: $rmRepoScoped$
stamps its own $myRepo$, $rmOrgScoped$ its own $myOrg$, so when the
inviter's $keyRepo$ (respectively $keyOrg$) differs the accept is
published to a subject the inviter does not subscribe to, and
$unsoundEmit$ latches.  Neither defect's mismatch is excluded by any
gate here: the inviter is visible but may live in another repository
\emph{and} another organization, which is exactly the case both defects
get wrong and pure-identity routing gets right.

%%----------------------------------------------------------------------
\subsection{MutualAutoAccept}
%%----------------------------------------------------------------------

Simultaneous mutual invites: two sessions each invited the other and
both are in $tpInviting$.  Without a tie-break, both would wait forever
for an accept.  The deterministic rule uses $keyBelow$: the
lexicographically-lower session key stays the inviter, and the higher
key auto-accepts.  This schema models the higher side --- an invite
from the very session we invited ($partnerKey?$) arrives, and our key
is above it ($partnerKey? \mapsto myKey \in keyBelow$), so we transition
straight to connected (publishing an accept as a side effect).  The
lower side takes no special operation; it simply receives our accept
via $DrainAcceptWhileInviting$.  Since $keyBelow$ is a total order,
exactly one side auto-accepts --- the handshake always completes.

\begin{schema}{MutualAutoAccept}
  \Delta TalkState \\
  partnerKey? : SESSIONKEY
\where
  phase = tpInviting \\
  partnerKey? \neq myKey \\
  partnerKey? = partnerKey \\
  (partnerKey? \mapsto myKey) \in keyBelow \\
  phase' = tpConnected \\
  myUser' = myUser \\
  myTty' = myTty \\
  myKey' = myKey \\
  myRepo' = myRepo \\
  myOrg' = myOrg \\
  partner' = partner \\
  partnerTty' = partnerTty \\
  partnerKey' = partnerKey \\
  pendingInvites' = pendingInvites \\
  queue' = queue \\
  queueLen' = queueLen \\
  visible' = visible \\
  routeMode' = routeMode \\
  ((unsoundEmit' = ztrue) \iff
    (unsoundEmit = ztrue \lor
      (routeMode = rmRepoScoped \land myRepo \neq keyRepo~partnerKey) \lor
      (routeMode = rmOrgScoped \land myOrg \neq keyOrg~partnerKey)))
\end{schema}

%%----------------------------------------------------------------------
\subsection{SendMessage}
%%----------------------------------------------------------------------

User types a line during talk mode.  Publishes a message
notification to the partner.

\begin{schema}{SendMessage}
  \Delta TalkState \\
  body? : MSGBODY
\where
  phase = tpConnected \\
  phase' = phase \\
  myUser' = myUser \\
  myTty' = myTty \\
  myKey' = myKey \\
  myRepo' = myRepo \\
  myOrg' = myOrg \\
  partner' = partner \\
  partnerTty' = partnerTty \\
  partnerKey' = partnerKey \\
  pendingInvites' = pendingInvites \\
  queue' = queue \\
  queueLen' = queueLen \\
  visible' = visible \\
  routeMode' = routeMode \\
  ((unsoundEmit' = ztrue) \iff
    (unsoundEmit = ztrue \lor
      (routeMode = rmRepoScoped \land myRepo \neq keyRepo~partnerKey) \lor
      (routeMode = rmOrgScoped \land myOrg \neq keyOrg~partnerKey)))
\end{schema}

Sending a message changes no talk component other than the routing
observation: the message is published via NATS core (fire-and-forget,
ephemeral), and $unsoundEmit$ latches only if an over-scoped defect
stamps it with the sender's own repository or organization rather than
the partner's identity.  Under $rmIdentity$ the message carries
$subjectOf~partnerKey = partnerKey$ and the flag is untouched --- so
messages, like accepts, reach a partner in any repository or
organization.  This is the conversation-side of the same routing
guarantee that governs the accept: the earlier draft's $\Xi TalkState$
under-modelled the send, treating a message as observationally inert,
which hid the very mis-routing today's defect exhibits.

%%----------------------------------------------------------------------
\subsection{LocalEnd}
%%----------------------------------------------------------------------

User types \texttt{end} during talk mode.  Publishes an end
notification and transitions to idle.

\begin{schema}{LocalEnd}
  \Delta TalkState
\where
  phase = tpConnected \\
  phase' = tpIdle \\
  partner' = myUser \\
  partnerTty' = myTty \\
  partnerKey' = myKey \\
  queue' = queue \\
  queueLen' = queueLen \\
  myUser' = myUser \\
  myTty' = myTty \\
  myKey' = myKey \\
  myRepo' = myRepo \\
  myOrg' = myOrg \\
  pendingInvites' = pendingInvites \\
  visible' = visible \\
  routeMode' = routeMode \\
  ((unsoundEmit' = ztrue) \iff
    (unsoundEmit = ztrue \lor
      (routeMode = rmRepoScoped \land myRepo \neq keyRepo~partnerKey) \lor
      (routeMode = rmOrgScoped \land myOrg \neq keyOrg~partnerKey)))
\end{schema}

%%----------------------------------------------------------------------
\subsection{CancelInvite}
%%----------------------------------------------------------------------

User types \texttt{end} while waiting for an accept (inviting
phase).  Cancels the invitation and returns to idle.

\begin{schema}{CancelInvite}
  \Delta TalkState
\where
  phase = tpInviting \\
  phase' = tpIdle \\
  partner' = myUser \\
  partnerTty' = myTty \\
  partnerKey' = myKey \\
  queue' = queue \\
  queueLen' = queueLen \\
  myUser' = myUser \\
  myTty' = myTty \\
  myKey' = myKey \\
  myRepo' = myRepo \\
  myOrg' = myOrg \\
  pendingInvites' = pendingInvites \\
  visible' = visible \\
  routeMode' = routeMode \\
  unsoundEmit' = unsoundEmit
\end{schema}

%%----------------------------------------------------------------------
\section{System Invariants}
%%----------------------------------------------------------------------

The following properties must hold in all reachable states:

\begin{enumerate}
\item \textbf{Idle implies no partner.}  When $phase = tpIdle$,
  $partner = myUser$ and $partnerTty = myTty$ (sentinel values).

\item \textbf{Queue bounded.}  $\# queue \leq maxQueueLen$.

\item \textbf{Self-echo never queued.}  For all $n$ in $\ran queue$,
  $n.nfromKey \neq myKey$.

\item \textbf{Pending invites only while idle.}  Invites are consumed
  on entering talk mode ($RespondToInvite$) or remain in the set
  until the next \texttt{talk} command.  The set is not modified
  during the inviting or connected phases (except by
  $RespondToInvite$ which transitions from idle).

\item \textbf{Handshake symmetry.}  If session~A is in $tpInviting$
  with $partner = B$, then session~B either is in $tpIdle$ (invite not
  yet seen), transitions to $tpConnected$ via $RespondToInvite$, or ---
  when B has also invited A --- resolves via $MutualAutoAccept$.  In the
  mutual case the $keyBelow$ tie-break guarantees exactly one side
  auto-accepts, so neither blocks.

\item \textbf{Mutual hangup.}  $LocalEnd$ on one side publishes
  $ntEnd$.  $DrainEnd$ on the other side transitions to $tpIdle$.
  After both operations, both sides are idle.

\item \textbf{Session-scoped delivery.}  A notification is queued only
  when $nto = myKey$.  A talk therefore reaches exactly the named
  session, never a user's other sessions ($ReceiveNotForSession$ drops
  the rest).

\item \textbf{Self-talk disallowed.}  $SendInvite$ requires
  $targetKey? \neq myKey$; a session can never open a talk to itself.

\item \textbf{Invite supersession.}  $pendingInvites$ is a function, so
  at most one pending invite exists per user; a newer invite overwrites
  the older via $\oplus$ in $DrainInvite$.

\item \textbf{Deterministic mutual accept.}  $keyBelow$ is a strict
  total order, so for two distinct session keys exactly one satisfies
  the $MutualAutoAccept$ guard.  Simultaneous mutual invites always
  resolve; the handshake never deadlocks.

\item \textbf{Accept consent boundary.}  An inviter transitions to
  $tpConnected$ only via $DrainAcceptWhileInviting$ or
  $MutualAutoAccept$, and both guard on $partnerKey$: the accept's
  origin ($(head~queue).nfromKey = partnerKey$) or the mutual invite's
  origin ($partnerKey? = partnerKey$) must equal the session we
  invited.  Since $SendInvite$ sets $partnerKey$ to $targetKey?$ and no
  operation reassigns it while inviting, a session in $tpConnected$
  reached through the invite path is connected to exactly the peer it
  invited.  An accept from any other origin is drained by
  $DrainForeignAccept$ without changing phase, so a third party who
  publishes an accept addressed to the inviter ($nto = myKey$) cannot
  forge the invited peer's consent.

  The same $partnerKey$ anchor governs the \emph{conversation}, not
  only the connect.  In $tpConnected$, $DrainMessage$ surfaces a
  message and $DrainEnd$ acts on an end only when
  $(head~queue).nfromKey = partnerKey$; a frame from any other origin
  is neither displayed in the partner's place nor able to hang up the
  call, but is drained by $DrainForeignMessage$ or $DrainForeignEnd$
  with no effect beyond dequeue.  A session in $tpConnected$ therefore
  both connects to, and converses with, exactly the peer it invited or
  accepted.

\item \textbf{Delivery soundness.}  $routeMode = rmIdentity \implies
  unsoundEmit = zfalse$.  Under identity routing no session ever
  publishes a reply to a subject other than its addressee's own, so
  every reply reaches the addressed \texttt{@user:tty} whatever
  organization \emph{or} repository either party is in.  This is the
  state invariant of the correct design; Section~\ref{sec:routing} sets
  it beside the two reachable strands --- one per over-scoped defect ---
  it excludes.

\item \textbf{Visibility is the only routing gate.}  $SendInvite$
  requires $targetKey? \in visible$ and $RespondToInvite$ requires the
  inviter $pendingInvites~target? \in visible$; there is no organization
  or repository precondition on either.  A frame enters the queue
  ($ReceiveNotification$, $ReceiveOverflow$) only when it arrives on this
  session's own subject $subjectOf~myKey = myKey$ --- the identity, and
  nothing else.  A session therefore addresses exactly the peers it can
  see, and each is reachable whatever its organization or repository; the
  invariant $phase \neq tpIdle \implies partnerKey \in visible$ records
  that a partner is always visible.  Neither $keyOrg$ nor $keyRepo$
  constrains the route: both are presence and visibility attributes, and
  a reply routes to $partnerKey$ from locally-held state, with nothing to
  resolve and nothing to persist.
\end{enumerate}

%%----------------------------------------------------------------------
\section{Routing Verification}
\label{sec:routing}
%%----------------------------------------------------------------------

The addressing layer is verified by four machine-checked claims, all at
the routing scope of Section~\ref{sec:addressing} (two session keys, two
repositories, two organizations, so the peer differs from the sender in
both repository \emph{and} organization at once).  They are stated so
that ProB decides each: the first as a state invariant, the rest as
reachability goals.  The correct discipline is $rmIdentity$ --- the
subject is the identity and a reply routes to $partnerKey$; the two
defects over-scope it --- $rmRepoScoped$ keys the subject on the
repository, $rmOrgScoped$ on the organization, and in each the responder
routes on its \emph{own} attribute.  Model checking the scope visits all
$2{,}089$ states with $invariant\_violated = 0$ and no deadlock, and
covers every operation but the two connected-phase forged-frame
operations (which need a third key, exhausted at the forgery scope).

\medskip\noindent\textbf{(R1) Delivery soundness (holds).}  As a state
invariant of $TalkState$,

\begin{verbatim}
routeMode = rmIdentity  =>  unsoundEmit = zfalse
\end{verbatim}

Under identity routing $unsoundEmit$ never latches: a reply is stamped
$subjectOf~partnerKey = partnerKey$, which is the subject the partner,
and only the partner, subscribes to.  Every reply is therefore delivered
to the addressed \texttt{@user:tty}, independently of which organization
\emph{or} repository either session runs in.  Nothing but the peer's
identity --- held locally, taken from the frame being answered --- is
read on the routing path.

\medskip\noindent\textbf{(R2a) The strand is reachable under
repo-keyed routing (holds).}  As a reachability goal,

\begin{verbatim}
probcli ... -goal "unsoundEmit=ztrue & routeMode=rmRepoScoped"
\end{verbatim}

ProB returns a counterexample.  With $keyRepo = \{k_1 \mapsto r_2,
k_2 \mapsto r_1\}$, $keyOrg = \{k_1 \mapsto o_2, k_2 \mapsto o_1\}$ and
$routeMode = rmRepoScoped$, a session whose key is $k_2$ ($myRepo = r_1$)
has $k_1$ in $visible$ and invites it; the two connect
($MutualAutoAccept$); but when it publishes its accept to the partner
$k_1$, the repo-keyed discipline stamps the frame with its \emph{own}
$myRepo = r_1$ rather than $keyRepo~k_1 = r_2$.  The partner's subject
is $(r_2, \ldots)$, so the accept is delivered to $ReceiveNotForSubject$
and dropped: the inviter never receives it and strands until its
transport TTL.  The peer was \emph{visible} --- routing was permitted ---
yet the repo coordinate in the subject broke delivery.  This is
biff-e9u.

\medskip\noindent\textbf{(R2b) The strand is reachable under org-keyed
routing (holds).}  As a reachability goal,

\begin{verbatim}
probcli ... -goal "unsoundEmit=ztrue & routeMode=rmOrgScoped"
\end{verbatim}

ProB returns a counterexample of the same shape with the organization in
the place of the repository: with $keyOrg = \{k_1 \mapsto o_1,
k_2 \mapsto o_2\}$ and $routeMode = rmOrgScoped$, a session $k_2$
($myOrg = o_2$) invites the visible $k_1$, they connect, and its accept
is stamped with its own $myOrg = o_2$ rather than $keyOrg~k_1 = o_1$;
the partner's subject is $(o_1, \ldots)$, so the frame is dropped and the
inviter strands.  Org-scoping over-scopes exactly as repo-scoping does:
a \emph{visible} peer in another organization is unreachable.  Both
defects are the same mistake --- a coordinate other than the identity in
the subject --- and both are exhibited here.

The two halves of ``never receives'' are each a theorem of this one
machine and compose across two instances, exactly as
Section~\ref{sec:introduction}'s two-session reading intends: the
\emph{sender} latches $unsoundEmit$ because it stamped a subject
$\neq subjectOf(\mathrm{addressee})$, and any \emph{receiver} whose
subject differs from a frame's arrival subject drops it by
$ReceiveNotForSubject$.

\medskip\noindent\textbf{(R3) Both strands are unreachable under
identity routing (holds).}  As a reachability goal,

\begin{verbatim}
probcli ... -goal "unsoundEmit=ztrue & routeMode=rmIdentity"
\end{verbatim}

ProB reports no such state after visiting all $2{,}089$: under
pure-identity routing neither strand configuration exists.  (R1) and
(R3) are the same fact seen twice --- as an invariant that holds and as
a bad state that is unreachable --- which is the precise sense in which
routing on the identity alone makes the strand \emph{impossible} rather
than merely \emph{unobserved}, for a peer differing in \emph{both}
organization and repository.

A full cross-attribute conversation is reachable at this scope as a
corollary of coverage: with $routeMode = rmIdentity$ and the two keys in
different repositories and organizations, $SendInvite$,
$RespondToInvite$, $DrainAcceptWhileInviting$, $SendMessage$,
$DrainMessage$, $LocalEnd$ and $DrainEnd$ are all covered, so a session
completes
invite\,$\rightarrow$\,accept\,$\rightarrow$\,message\,$\rightarrow$\,end
with every frame sound.  By the symmetry of the two instances the peer
does likewise, so two visible sessions in different repositories
\emph{and} different organizations complete a full talk with every frame
reaching the correct identity.  Consent is untouched by any of this: the
session-key guards ($DrainAcceptWhileInviting$, $DrainMessage$,
$DrainEnd$ and their $DrainForeign$ counterparts) decide \emph{whether} a
delivered frame is honoured, orthogonally to the subject that decides
\emph{where} it goes (DES-046, verified at the forgery scope).

%%----------------------------------------------------------------------
\section{Ownership: this specification and the session model}
\label{sec:ownership}
%%----------------------------------------------------------------------

Addressing spans two specifications, and the split is deliberate.

\begin{itemize}
\item $session\-model.tex$ \textbf{owns the entities and the
  resolution.}  Organizations, repositories, members, machines, clones,
  sessions and the session attributes --- in particular $srepo :
  SESSION \pfun (NAME \cross NAME)$, a session's (org, repo) --- live
  there, together with the referential invariants that keep them
  consistent.  It also owns what it currently defers: the resolution of
  the \texttt{@user:tty} address to a session (its ``Addressing
  syntax'' extension point) and the visibility relation the
  \texttt{peers} and \texttt{orgs} configuration induces.  In this split
  both $srepo$ \emph{and} the organization it names are \emph{visibility
  and presence} attributes --- they say where a session appears in
  \texttt{who}, \texttt{finger} and discovery, and which peers each
  session can resolve --- and they are nothing more.

\item $talk.tex$ \textbf{owns the routing discipline.}  It takes the
  resolution and the visibility set as given, through $keyUser$,
  $subjectOf$ and $visible$ (Section~\ref{sec:addressing}), and states
  the one fact the transport must honour: a reply's subject is
  $subjectOf$ of the addressed identity, which is that identity itself
  --- computed by the sender from the peer's identity alone, never from
  the organization or repository of either party.  The single gate on
  the route is visibility ($targetKey? \in visible$); the consent and
  footprint guarantees (the session-key guards, the queue bound) are also
  its own.  $keyRepo$ and $keyOrg$ appear here only to state that
  repository and organization are visibility attributes and to give the
  two modelled defects a coordinate to mis-route on.
\end{itemize}

The bridge between the two is DES-035's global uniqueness: a
\texttt{@user:tty} names exactly one session across all repositories and
organizations, so $SESSIONKEY$ here is $session\-model.tex$'s $SESSION$
addressed by its \texttt{@user:tty}, and $keyRepo$, $keyOrg$ are $srepo$
and its organization projected through that address.  We do not restate
$srepo$; to do so would be to keep two copies of one fact, and the
copies would drift.

One reconciliation is owed to $session\-model.tex$, and it is a
recommendation rather than a change made here.  That specification lists
``the \texttt{@user:tty} \ldots resolution rules'' as a natural
extension it does not yet cover, and it does not yet model visibility.
To make $keyUser$, $keyRepo$ and $keyOrg$ derivations rather than fresh
axioms, $session\-model.tex$ should expose an address function ---
$saddr : SESSION \inj (USER \cross TTY)$, injective by DES-035 --- from
which they follow by projection of $srepo$; and it should expose the
visibility relation $visible$ derives from.  Until it does, the resolver
axioms and the $visible$ set of Section~\ref{sec:addressing} are
$talk.tex$'s local statement of that interface, consistent with $srepo$
by construction.  What routes is the identity, and the only gate is
visibility; organization and repository are presence attributes on both
sides.  Where a would-be design choice conflicts with this --- namely
keying the subject on the repository \emph{or} the organization and
routing a reply on the \emph{sender's} attribute --- the architecture
(a globally-unique \texttt{@user:tty} that is reachable wherever it is
visible) is authoritative, and the model proves each such choice is one
that strands a visible peer (R2a, R2b) while identity routing does not
(R1, R3).

%%----------------------------------------------------------------------
\section{Threat Model: forgeable frames and the withdrawal primitive}
\label{sec:threat}
%%----------------------------------------------------------------------

Talk inherits the BSD trust model of the programs it is named for: a frame
is trusted for what it claims to be.  Notifications travel on a per-identity
NATS subject and carry their own $nfrom$, $nfromTty$ and $nfromKey$;
nothing in the protocol binds those fields to a verified sender, and a frame
from any key --- visible or not --- may be published to our subject.  A frame
is therefore \emph{forgeable}, and the specification's guarantees are stated
against that fact rather than against an assumption of authenticity.

The consent boundary (invariant~11) is robust under forgery on the
\emph{connect} side.  An accept is honoured only when
$(head~queue).nfromKey = partnerKey$, so a forged accept connects the inviter
to nobody it did not itself invite: to be honoured, the forgery must carry
the exact ephemeral session key the inviter chose as its target, which a
third party does not hold.  Forging a session identity, not merely a
username, is the bar.

The same boundary now governs the \emph{conversation}, not only the
connect, and this is a correction: earlier drafts guarded the
handshake but left the connected phase open, so the guarantee below
overclaimed relative to the model.  A message is surfaced, and an end
acts, only when $(head~queue).nfromKey = partnerKey$ ($DrainMessage$,
$DrainEnd$); a frame from any other origin is drained without effect
by $DrainForeignMessage$ or $DrainForeignEnd$.  A forged $ntMessage$
carrying $nto = myKey$ therefore cannot be displayed in the partner's
place, and a forged $ntEnd$ cannot hang up the call: to be honoured,
either forgery must carry the exact ephemeral session key of the
connected peer, which a third party does not hold.  The guarantee that
no forged frame can \emph{speak in another's place} is now true by
construction, and its residual reduces to partner-key forgery --- the
same session-identity bar the consent boundary already sets, no higher
and no lower.

A withdrawal frame---the \emph{ntWithdraw} of the sister specification
$notification.tex$, which $talk.tex$ would model as a fifth $NotifType$---sits
on the \emph{availability} side of this boundary, and the honest reading is
that it introduces a targeted invite-suppression primitive.  A withdrawal
removes a pending invite; a party who can forge one can cancel an invitation
the recipient has not yet answered, denying the handshake.  This is a
denial-of-availability, not a breach of consent: withdrawal only \emph{removes}
pending state and never drives a transition into $tpConnected$, so no forged
frame can connect a session to a peer or speak in another's place.  The
consent boundary is untouched.

The residual, then, is availability, and the session-key guard on the
recipient's side---$notification.tex$'s $WithdrawArrive$, which applies a
withdrawal only when $talkPending~inviter? = withdrawKey?$ and otherwise drops
it via $WithdrawForeign$---raises the bar for that residual to the same height
as the accept guard.  Before the guard, a withdrawal matched on the inviter's
username alone: any party who knew the username could suppress the invite, and
even an \emph{honest} but stale withdrawal from the inviter's own earlier
session could delete a live invite from a later one, since the transport gives
no cross-session ordering.  After the guard, a suppressing withdrawal must
carry the inviter's ephemeral session key, exactly as a connecting accept
must---so targeted suppression now demands the same session-identity forgery
that consent already demands, and the non-adversarial reordering is closed
outright.  The guard does not make withdrawal unforgeable; it makes a forged
or stale withdrawal no cheaper than a forged accept, which is the most the
BSD-ephemeral trust model admits.  The corresponding obligation on the
recipient is the $WithdrawArrive$ key match of $notification.tex$.

\end{document}
